\section{Implications for the Agentic Web}
\label{sec:vision}

The Brave New Ideas track asks what the Web could become if a capability existed that does not exist today. We argue that action-centered execution auditing is that kind of capability, and we sketch what changes if it works at Web scale.

\paragraph{From content auditing to execution auditing.}
The Web's existing safety infrastructure---spam filters, content moderation, phishing detection, provenance signals---inspects \emph{content} at fetch time. For agents, the dangerous event is not the fetched byte but the executed operation, and attacks can be invisible in content and outcome alike~\cite{lee2026will}. If browsers, agent runtimes, and tool protocols emit trajectory graphs as a first-class artifact---the same way access logs and network traces are emitted today---then every agent-mediated action carries an auditable record of what state motivated it, what reasoning justified it, and what it changed. Detection and localization over such records generalize beyond prompt injection: any policy violation, error, or fraud that manifests in execution structure becomes addressable with the same substrate, including failures that outcome-based evaluation misses~\cite{mohammadi2026trajectoryjudge,deng2026defa}. This is a concrete, engineering-shaped path from today's ad-hoc guardrails toward execution-layer trust for the agentic Web~\cite{yang2025agentic,ng2026agent}.

\paragraph{Experiences as a portable safety layer.}
The \sexp{} format is deliberately an artifact, not a log line: it is small, structured, grounded in cited node identifiers, and written for a consumer agent rather than a human analyst. If experiences accumulate across runs, sites, and agent configurations, safety knowledge stops being bound to a single deployment---the failure a search agent suffers today can protect a booking agent tomorrow, and the memory-and-skill externalization trend of the field~\cite{zhou2026externalization,zhao2023expel,allard2026experiential} extends from competence to safety. The defense conditions of \S\ref{sec:defense} are designed to measure exactly the property this vision needs: whether an experience grounded in localized evidence transfers beyond the session that produced it.

\paragraph{Open problems the substrate raises.}
An execution-audit layer creates problems we do not solve here. \emph{Standardization}: trajectory-graph emission needs schema-level agreement among runtimes, in the spirit of emerging tool-protocol standards~\cite{hou2025model,zhang2025security}. \emph{Privacy}: graphs record how agents act on behalf of users, so emission policies must decide what is stored, for how long, and under whose access control. \emph{Experience pollution}: shared experience pools can be poisoned, and an attacker who understands the format can craft content that generates misleading experiences; provenance and robustness of the experience channel are prerequisites for the ecosystem, echoing lessons from repository- and knowledge-poisoning research. \emph{Governance}: if audits attribute actions to specific agent configurations, the substrate becomes evidence in accountability disputes, which is a feature for site operators and a liability question for agent vendors. We state these as open problems because they require the community-level agreement that a research paper cannot produce; what this paper contributes is the demonstration that the substrate's core operation---joint detection, operation-level localization, and grounded experience generation---is learnable at all.
